\section{The problem}\label{ch:problem}

\paragraph{Federated learning} In federated learning, data stay on the devices or institutions that own them, the
\emph{clients}. A \emph{server} coordinates training. In each \emph{round} the server sends a model to the clients,
each client trains it for a while on its own data, and the server combines the returned models into a new one. The
standard way of combining is FedAvg: the server averages the clients' parameters, coordinate by coordinate. No raw
data leave a client; only parameters or other summaries do.

\paragraph{Different models} FedAvg needs every client to train the same network, because only then do the $j$-th
parameters of two clients mean the same thing. In practice clients run different hardware, so they may run different
networks: a phone runs a small convolutional network, a hospital server a larger one. This is
\emph{model-heterogeneous} federated learning. Parameters of different architectures cannot be averaged, so the
clients must share knowledge in another form. The paper studies one such form, a shared generator of synthetic
features (\cref{ch:gefl}).

\paragraph{Skewed and long-tailed labels} The second difficulty is in the labels. Two properties combine.
\begin{itemize}
\item \emph{Label skew}: clients hold different classes in different proportions. A client may have hundreds of images
of one class and none of another.
\item \emph{A global long tail}: over the whole federation, a few classes are frequent and many are rare.
\end{itemize}
The paper builds both in a controlled way. The class sizes follow
\[
n_c\propto\rho^{-c/(C-1)},\qquad c=0,\dots,C-1,
\]
so class $0$ is the largest, class $C-1$ the smallest, and the \emph{imbalance ratio} $\rho=n_0/n_{C-1}$ fixes how
steep the tail is. The paper uses $\rho=100$. Each class is then divided among the $K$ clients by a Dirichlet draw: for
class $c$ a vector of shares $(p_{1c},\dots,p_{Kc})$ is drawn from $\mathrm{Dirichlet}(\alpha,\dots,\alpha)$, and client
$k$ receives a fraction $p_{kc}$ of the class. A small concentration $\alpha$ gives very uneven shares; the paper uses
$\alpha=0.5$.

\paragraph{The running example} The text returns to one concrete case throughout: CIFAR-10-LT with $\rho=100$, $K=10$
clients and $\alpha=0.5$, split with seed 3. The run logs give its class sizes,
\[
(n_0,\dots,n_9)=(4800,\,4800,\,3620,\,2170,\,1301,\,779,\,467,\,280,\,168,\,100),
\]
18,485 training images in all. The profile asks for more than 4,800 images of classes 0 and 1, but CIFAR-10 has only
5,000 per class and 200 are held out for validation, so these two classes are capped. The realised ratio is
therefore $4800/100=48$, not 100. The paper still names the setting by its target ratio $\rho=100$ and gives the
realised ratio in its dataset table (\cref{tab:s-datasets}).

The rare class 9 shows what the split does to a rare class. Of its 100 images, client 7 holds 63. The other holders
have between 1 and 14 each, and clients 2 and 5 have none. (Clients are numbered 1 to 10 here; in the run logs they
are 0 to 9.) A class that is rare in the federation is held by few clients, and most of them see it only a few
times.

\paragraph{How the loss is measured} Plain accuracy hides what happens to rare classes. Take a model that recognises
classes 0--4 with recall 90\% and classes 5--9 with recall 10\%. On a test set with the training proportions above,
its accuracy is
\begin{align*}
&\frac{0.9\cdot(4800+4800+3620+2170+1301)+0.1\cdot(779+467+280+168+100)}{18485}\\
&\qquad=\frac{15021.9+179.4}{18485}\approx82\%,
\end{align*}
because the first five classes hold 90\% of the images. Yet the model fails on half of the classes. \emph{Balanced
accuracy} is the mean of the per-class recalls, here $(5\cdot0.9+5\cdot0.1)/10=50\%$: every class counts equally,
whatever its size. It is the paper's primary metric. The test sets of CIFAR-10 and CIFAR-100 are balanced (1,000 and
100 images per class), so there the two measures coincide; on the naturally imbalanced datasets they do not.

With ten client models, one per architecture, the paper reports the mean over the ten models of each model's balanced
accuracy,
\[
\text{balanced accuracy}=\frac1G\sum_{g=1}^{G}\frac1C\sum_{c=1}^{C}r_{gc},
\]
where $r_{gc}$ is the recall of class $c$ by the model of architecture $g$ and $G=10$ is the number of architectures.

\begin{table}[t]
\centering\small
\caption{The datasets as used in the paper (the paper's dataset table): number of classes, training images held by
the clients, test images, and the realised ratio between the largest and the smallest class. The first two are the
long-tailed benchmarks built with $\rho=100$; the others keep their own class sizes.}
\label{tab:s-datasets}
\slottab{datasets}{}
\end{table}

\paragraph{Two protocols} The paper tests its method in two settings.
\begin{itemize}
\item \emph{Different client models}, the setting the method is built for: ten clients, each with its own head
architecture, on the long-tailed benchmarks CIFAR-10-LT and CIFAR-100-LT and on six datasets whose classes are
naturally imbalanced (\cref{tab:s-datasets}).
\item \emph{One shared model}, the standard protocol of federated long-tail research: twenty clients train one
ResNet-8 together. Here the specialised long-tail methods of the literature apply, so this setting shows how the
method compares with them outside its own design.
\end{itemize}

The question the rest of this text builds on is how clients that run different models can share knowledge at all.
